\scanpage{13}
que $A'$ est muni d'une \uline{donnée de descente} relativement à $S'\mathbin{-}S$. En particulier, \deleted{$\unsure{\ell}$} ${}_\ell A'$ est muni d'une telle donnée de descente; \deleted{\unsure{que}} \deleted{\unsure{et}} Or $S'\mathbin{-}S$ est \inserted{un morphisme} de descente effective pour les schémas relatifs finis étales (SGA 1 VIII), donc on trouve un \deleted{\unsure{tel}} tel rev. \deleted{$\unsure{R}$} \inserted{en groupes} $N$ de $S$, tel que $N\mid U\simeq{}_{\ell}(A_U)$. Cela montre que l'on peut trouver un morphisme \inserted{\uline{fini}} \uline{étale} surjectif $S_1\to S$ qui \unsure{splitte} ${}_\ell(A_U)$, \deleted{Remplaçant alors $S'$} et on peut supposer $S'$ \deleted{fini} \uline{étale} sur $S$, donc plat. Alors ($A'$ est projectif sur $S'$) on \deleted{trouve une donnée de \unsure{des}} \inserted{peut \uline{descendre}} $A'$ en $A$, qui résoud le pb.

\margin{\fbox{\parbox{.8\linewidth}{\deleted{le normalisé} tel qu'il existe $S'\mathbin{-}S$ fini surjectif, avec $S'$ géom unib.}}}
d) Cas $S$ réduit, \unsure{(}$S'$ fini sur $S$. Utilisant le fait que $S$ est géom unibranche en les points de $S-U$, on voit que l'image inverse $U''=U'\times_U U'$ de $U$ dans $S''=S'\times_S S'$ est dense dans $S''$.
\margin{\diagram{\begin{tikzcd}[column sep=small,row sep=small]
U\arrow[d,phantom,"\rotatebox{90}{$\unsure{\subset}$}"] & U'\arrow[l]\arrow[d,phantom,"\rotatebox{90}{$\unsure{\subset}$}"] & U''\arrow[l,shift left=.6ex]\arrow[l,shift right=.6ex]\arrow[d,phantom,"\rotatebox{90}{$\unsure{\subset}$}"]\\
S & S'\arrow[l] & S''\arrow[l,equal]
\end{tikzcd}}}
\deleted{Par le cas b), on}
Par le cas c), on trouve sur $S'$ un schéma abélien $A'$, et par la pleine fidélité (pour $S'',U''$), une donnée de descente sur $A'$, relativement à $S'\to S$. Notons que le morphisme $S'\mathbin{-}S$, étant un épimorphisme fini, est un morphisme de descente pour les schémas abéliens, \deleted{\unsure{et}} et aussi pour les grpes $p$-divisibles. Donc on est ramené à
\margin{lemme à dégager \figure{\begin{tikzpicture}\draw (0,0)--(0,-.8);\end{tikzpicture}}}
\mspage{12}
